\documentclass[11pt]{article}
\usepackage[utf8]{inputenc}
\usepackage{amsmath,amssymb}
\usepackage[margin=1in]{geometry}
\usepackage{hyperref}
\emergencystretch=3em

\title{Joint limit of the normalised chart integrals for deterministic amplitudes}
\author{Claude, with Daniel Murfet}
\date{2026-09-07}

\begin{document}
\maketitle
\sloppy

\begin{abstract}
Astra~\#13 unit u158. In §4.3 only the phase $\xi_n$ is random; the amplitudes
$\eta_\varphi=\varphi\cdot\mathrm{prior}\cdot\mathrm{Jac}$ and $\eta_1=\mathrm{prior}\cdot\mathrm{Jac}$
are deterministic. We model this by \texttt{withY}~$y$, the map on the coefficient-pair space that
keeps the phase data and replaces the amplitude data by a fixed array $y$ ($1$-Lipschitz). For equal
starting exponents $p$ the leading exponent has no lower terms, and the pair of normalised chart
integrals $(Z_{N_n}[\varphi]/s_n,\,Z_{N_n}[1]/s_n)$ with $s_n=N_n^{-p}\log N_n$ converges jointly in
distribution to $(A_p(X;y_\varphi),A_p(X;y_1))$ whenever the random phase data converge in
distribution (\texttt{tendstoInDistribution\_normA\_pair}). Zero \texttt{sorry}/\texttt{axiom}.
\end{abstract}

\section{The things formalised}
{\small
\begin{verbatim}
def withY (ρ) (hρ) (y) (hy : WSummable ρ y) (a : CoeffPair) : CoeffPair :=
  ofPair ρ hρ (toX ρ a) y (wsummable_toX ρ hρ a) hy
theorem toX_withY / toY_withY / withY_apply_inl / withY_apply_inr
theorem norm_withY_sub_le : ‖withY … a - withY … a'‖ ≤ ‖a - a'‖
theorem lipschitzWith_withY / continuous_withY / measurable_withY
theorem lowerPart_leading_eq_zero (h₁ h₂ k₁ k₂) (hk₁ hk₂) (p T) (hp₁ hp₂) (A B) (N) :
    lowerPart (polesBelowGen h₁ h₂ k₁ k₂ p p T) A B p N = 0
theorem tendstoInMeasure_prodMk_zero (μ) (f g) (hf : TendstoInMeasure μ f l (fun _ => 0))
    (hg) :
    TendstoInMeasure μ (fun n ω => (f n ω, g n ω)) l (fun _ => (0, 0))
theorem tendstoInMeasure_normA_withY_sub ... (hpT : p < 2 * T) (y) (hy) (X) (Z) (hX)
    (Nseq) (hN) :
    TendstoInMeasure μ (fun n ω => normA … p (Nseq n) (withY … y hy (X n ω))
      - coeffA … p (withY … y hy (X n ω))) l (fun _ => 0)
theorem tendstoInDistribution_normA_pair ... (yφ y₁) (hyφ hy₁) (X) (hXm) (Z) (hX)
    (Nseq) (hN) :
    TendstoInDistribution (fun n ω => (normA … p (Nseq n) (withY … yφ … (X n ω)),
                                        normA … p (Nseq n) (withY … y₁ … (X n ω)))) l
      (fun ω => (coeffA … p (withY … yφ … (Z ω)), coeffA … p (withY … y₁ … (Z ω))))
      (fun _ => μ) μ'
\end{verbatim}
}

\section{Mathematics}
$\texttt{withY}\,y\,a$ has the same $\mathrm{inl}$-coordinates as $a$ and fixed
$\mathrm{inr}$-coordinates, so $\|\texttt{withY}\,y\,a-\texttt{withY}\,y\,a'\|\le\|a-a'\|$ termwise.
The leading exponent $p=\min\Lambda$ has empty lower part. The two normalised chart integrals are
each the corresponding coefficient plus an error that tends to zero in probability (the argument of
units~152--153 applied to the pushed-forward data $\texttt{withY}\,y\circ X_n$, which converge in
distribution by the continuous mapping theorem and are therefore norm-bounded in probability); the
pair of errors tends to zero in probability, and the pair of coefficients converges jointly by the
continuous mapping theorem applied to the continuous map
$a\mapsto(A_p(\texttt{withY}\,y_\varphi\,a),A_p(\texttt{withY}\,y_1\,a))$.

\section{Paper-facing meaning}
This is the joint numerator/denominator convergence that cor:empirical\_expectation's division
needs, in the paper's own setting (deterministic amplitudes, random phase). With the quotient
theorem of unit~157 and the positivity of unit~154 the chart posterior limit follows in the next
unit, including Astra~\#13's observation that the common factor $\int s^{p-1}e^{-\beta s^2+\beta sx_{00}}ds$
cancels, so that $Z_n[\varphi]/Z_n[1]\Rightarrow y_{\varphi,00}/y_{1,00}$.

\section{Lean notes}
\begin{itemize}
\item Coordinates of \texttt{ofPair} are exposed by \texttt{change} (the underlying function is a
\texttt{Sum.elim}); \texttt{lp.coeFn\_sub} then \texttt{Pi.sub\_apply} for differences.
\item \texttt{LipschitzWith.of\_dist\_le\_mul} with \texttt{NNReal.coe\_one}.
\item \texttt{Prod.mk\_sub\_mk} and \texttt{Prod.norm\_def} ($\|(a,b)\|=\max(\|a\|,\|b\|)$) reduce the
pair convergence in measure to the components via \texttt{measure\_union\_le}.
\item For \texttt{ENNReal}, \texttt{zero\_le} takes its argument implicitly; unused hypotheses are
underscored and the unused section variable is \texttt{omit}ted.
\end{itemize}

\end{document}
